\documentclass[10pt,a4paper]{amsart}
\usepackage[margin=19mm]{geometry}
\usepackage{amsmath,amssymb,mathtools}
\usepackage[scr=rsfs,cal=euler]{mathalfa}
\usepackage{xfrac}
\usepackage[unicode,hidelinks,bookmarks=false]{hyperref}
\newcommand{\ie}{i.e.\ }
\newcommand{\cf}{cf.\ }
\newcommand{\resp}{respectively\ }
\newcommand{\Subsection}{\subsection}
\providecommand{\nobreakdash}{\nobreak\hbox{-}}
\setlength{\emergencystretch}{2em}
\multlinegap0pt
\usepackage{bm}
\usepackage{bbm}
\usepackage{dsfont}
\usepackage{xspace}
\usepackage{cancel}
\usepackage{mathtools}
\usepackage{amsthm}
\makeatletter
\@ifundefined{thm}{\newtheorem{thm}{Thm}}{}
\@ifundefined{lemma}{\newtheorem{lemma}[thm]{Lemma}}{}
\makeatother
\newcommand{\mc}[1]{\mathcal{#1}}
\renewcommand{\bar}[1]{\overline{#1}}
\title[Classifying the complexity of models of arithmetic]{Classifying the complexity of models of arithmetic}
\author{David Gonzalez, Mateusz {\L}e{\l}yk, Dino Rossegger,, Patryk Szlufik}
\date{Source: arXiv:2507.12025v1 (CC BY 4.0)}
\begin{document}
\maketitle

\noindent\emph{Source and licence.} Classifying the complexity of models of arithmetic. Version arXiv:2507.12025v1 (CC BY 4.0), distributed under the Creative Commons Attribution 4.0 International licence, as recorded in the arXiv licence metadata for this exact version and shown on the abstract page https://arxiv.org/abs/2507.12025v1. Full rights notice: licenses/arxiv-2507.12025-CC-BY-4.0.txt.\par\smallskip

\noindent\textit{Editorial scope.} Counted unit: the Lemma whose source label is \texttt{lem:colorHelper} in the article (source0.tex lines 607--614, source label \texttt{lem:colorHelper}), delivered together with the article's own complete original proof of it (source0.tex lines 616--631, one \texttt{proof} environment of 16 lines). No mathematics is written, repaired or re-derived: statement and proof are the article's own text line for line. Required context retained uncounted: none. Every definition, display and label of those lines is retained unchanged. Omitted from this package and not redistributed: 1--606, 615--615, 632--1324. Nothing inside a retained excerpt is omitted or altered: every retained body is delivered exactly as the article's lines are written, so no omission receipt of any kind accompanies this package. Citations: the retained excerpts carry no citation command at all, so no bibliography entry of the article is transcribed and no reference is redistributed. Reference adaptation: none was needed and none was made; every reference inside the delivered unit resolves inside it, so no unresolved reference or citation is produced. Printed slips kept literal: no printed word, symbol, hypothesis or number of the counted statement or of its proof was altered. Layout and class: the article's own class is \texttt{amsart} with packages including amssymb, inputenc, hyperref, xcolor, pdfsync, graphicx, eucal, diagbox, tabularray, lmodern, cleveref, amsmath, amsthm, mathtools; the wrapper is the lane's pinned \texttt{amsart} editorial wrapper at 10pt on A4, which restates the theorem-like environments the retained text needs and adds the article's own macro definitions that the retained text needs (\texttt{\textbackslash bar}, \texttt{\textbackslash mc}), with extra wrapper packages (bm, bbm, dsfont, xspace, cancel, mathtools). The delivered numbering is the unit's own. Assets: no retained span contains a figure environment, a graphics-inclusion command, a PDF-page inclusion, a table or a TikZ picture; no image file is copied into this package, the package declares no asset dependency and no image file is redistributed with it. Licence: version arXiv:2507.12025v1 of this article is distributed under the Creative Commons Attribution 4.0 International licence, as recorded in the arXiv licence metadata for this exact version and shown on the abstract page https://arxiv.org/abs/2507.12025v1. A full rights notice is delivered at licenses/arxiv-2507.12025-CC-BY-4.0.txt. Characters: every retained excerpt is pure ASCII, so the delivered engine, XeLaTeX with its default Latin Modern text font, reports no missing character.
\begin{lemma}\label{lem:colorHelper}
    Let $LO_k$ be the theory of linear ordering with $2\leq k\leq\aleph_0$ colors. Let $\mc{L},\mc{K}\models LO_k$. Let $\bar{a}=a_1<\cdots<a_n\in\mc{L}$ and $\bar{b}=b_1<\cdots<b_n\in\mc{K}$ and by convention let $a_0=b_0=-\infty$ and $a_{n+1}=b_{n+1}=\infty$.
    \begin{enumerate}
        \item $(\mc{L},\bar{a})\leq_\alpha (\mc{K},\bar{b})$ if and only if for each $i\leq n$, $a_i$ has the same color as $b_i$ and $(a_i,a_{i+1})_{\mc{L}}\leq_\alpha (b_i,b_{i+1})_{\mc{K}}$.
        \item $\mc{L}\leq_1\mc{K}$ if and only if every finite, increasing sequence of colors in $\mc{K}$ is also present in $\mc{L}$.
    \end{enumerate}
    
\end{lemma}

\begin{proof}
We begin by showing the forward direction of Item (1).
Say that for some $i$, $(a_i,a_{i+1})_{\mc{L}}\not\leq_\alpha (b_i,b_{i+1})_{\mc{K}}$ witnessed by some formula $\psi$.
Relativizing the quantifiers in $\psi$ to being between the $i^{th}$ and $(i+1)^{st}$ elements of the tuple of parameters directly gives a $\hat{\psi}$ that witnesses $(\mc{L},\bar{a})\not\leq_{\alpha} (\mc{K},\bar{b})$.
Also, if $a_i$ and $b_i$ are different colors, $(\mc{L},\bar{a})\not\leq_0 (\mc{K},\bar{b})$ follows immediately.

The backward direction of Item (1) is a transfinite induction with the only interesting case being the successor case.
Say that $a_i$ has the same color as $b_i$ and $(a_i,a_{i+1})_{\mc{L}}\leq_{\beta+1} (b_i,b_{i+1})_{\mc{K}}$ for all $i\leq n$.
Consider $\bar d$ with $\bar d_i\subseteq \bar d$ in $(b_i,b_{i+1})$. Then there is $\bar c_i\in (a_i,a_{i+1})$ such that $(b_i,b_{i+1},\bar d_i)\leq_\beta ( a_i, a_{i+1},\bar c_i)$. 
For $i\leq n$ let $\bar{c}_i$ be the elements of $\bar{c}$ in $(b_i,b_{i+1})$.
Writing $\bar{a},\bar{c}=\bigcup \bar c_i$ as an increasing sequence $\bar{p}$ and $\bar{b},\bar{c}=\bigcup \bar d_i$ as an increasing sequence $\bar{q}$, the induction hypothesis gives that
\[(\mc{L},\bar{a},\bar{c})=(\mc{L},\bar{p})\geq_\beta(\mc{K},\bar{q})=(\mc{K},\bar{b},\bar{d}),\]
as desired.

Item (2) follows directly by observing that the quantifier-free type of a tuple is entirely determined by its ordering and the colors of the elements.
\end{proof}

\end{document}
